\documentclass[a4paper]{article}

% ============================================================
%  ENGINE -- LAYOUT & TYPESETTING LOGIC ONLY
%  Nothing below this block depends on the content above it.
%  An LLM filling in CV content should never
%  edit anything between here and "END ENGINE" -- only the
%  CONTENT ZONE that follows it. See the command reference at
%  the end of this block for the full editable API.
% ============================================================

% -------------------- 1. PACKAGES ---------------------------
% Page layout & structure
\usepackage[top=0.55in, bottom=0.55in, left=0.55in, right=0.55in]{geometry}
\usepackage{multicol}
\usepackage{fancyhdr}
% Lists & section titles
\usepackage{enumitem}
\usepackage{titlesec}
% Typography & symbols
\usepackage{charter}
\usepackage{fontawesome5}
\usepackage{latexsym}
\usepackage{marvosym}
\usepackage{pifont}
% Colors & hyperlinks
\usepackage{xcolor}
\usepackage{hyperref}

% -------------------- 2. PAGE LAYOUT -------------------------
\setlength{\footskip}{3.60004pt}
\urlstyle{rm}
\raggedbottom
\raggedright
\setlength{\tabcolsep}{0in}
\pagestyle{fancy}
\fancyhf{}
\renewcommand{\headrulewidth}{0pt}
\renewcommand{\footrulewidth}{0pt}

% -------------------- 3. COLORS & HYPERLINKS ------------------
\hypersetup{
  colorlinks,
  linkcolor={red!50!black},
  citecolor={blue!50!black},
  urlcolor={blue!50!black}
}

% -------------------- 4. SECTION TITLE FORMAT -----------------
\titleformat{\section}{\vspace{-10pt}\bfseries\raggedright\large}{}{0em}{}[\color{black}\titlerule \vspace{-6pt}]

% -------------------- 5. LIST CONFIGURATION --------------------
% Level 1: main entries (roles, projects, education)
\setlist[itemize,1]{leftmargin=*, topsep=6pt, itemsep=6pt, parsep=0pt, before={\vspace{2pt}}}
% Level 2: bullets within an entry
\setlist[itemize,2]{leftmargin=2em, topsep=3pt, itemsep=2pt, parsep=0pt}
% Compact list -- used by Technical Capabilities & Languages
\newlist{compactitemize}{itemize}{1}
\setlist[compactitemize]{label=\textbullet, labelsep=0.5em, leftmargin=*, topsep=4pt, itemsep=2pt, parsep=0pt, before={\vspace{2pt}}}
\renewcommand{\labelitemii}{$\circ$}
% Dash sub-bullets -- one level deeper than a normal entry (e.g. inside Training & Development)
\newcommand{\resumeSubBulletListStart}{\begin{itemize}[label=--, nosep, leftmargin=2em]}
\newcommand{\resumeSubBulletListEnd}{\end{itemize}}

% -------------------- 6. HEADER BUILDING BLOCKS -----------------
% Vertical-bar separator between contact items
\newsavebox\ContactSeparatorBox
\sbox\ContactSeparatorBox{$|$}
\newcommand{\ContactSeparator}{\unskip\cleaders\copy\ContactSeparatorBox\hskip\wd\ContactSeparatorBox\ignorespaces}
% Centered, slightly-leaded header block
\newenvironment{documentHeader}{\setlength{\topsep}{0pt}\par\kern\topsep\centering\linespread{1.5}}{\par\kern\topsep}

% -------------------- 7. ENTRY COMMANDS --------------------------
% A single bullet, fixed 10/12 font sizing for consistent leading
\newcommand{\resumeItem}[1]{\item{\fontsize{10}{12}\selectfont #1}}
% A bullet with a bold leading label (Technical Capabilities, Languages)
\newcommand{\resumeSubItem}[2]{\resumeItem{\textbf{#1} #2}}
% Two-line entry: Org/Institution + Location, then Role + Date
% Used by: Training & Development, Model Releases (Experience)
\newcommand{\resumeSubheading}[4]{%
  \item
  \begin{tabular*}{0.97\textwidth}{l@{\extracolsep{\fill}}r}
    \textbf{#1} & #2 \\
    \textit{#3} & \textit{#4}
  \end{tabular*}%
}
% One-line entry: Title + italic subtitle, Date
% Used by: Projects & Research, Industry Initiatives
\newcommand{\projectSubheading}[3]{%
  \item
  \begin{tabular*}{0.97\textwidth}{l@{\extracolsep{\fill}}r}
    \textbf{#1} \textit{#2} & #3 \\
  \end{tabular*}%
}
% Two-line entry purpose-built for certifications / courses / evaluations:
% Title + Date on row 1, issuing body in italics on row 2.
% (This replaces reusing \resumeSubheading with a blank 2nd argument,
%  which leaves a stray empty cell next to the title.)
\newcommand{\cvCertification}[3]{%
  \item
  \begin{tabular*}{0.97\textwidth}{l@{\extracolsep{\fill}}r}
    \textbf{#1} & #3 \\
    \textit{#2} & \\
  \end{tabular*}%
}

% -------------------- 8. LIST WRAPPERS -----------------------------
\newcommand{\resumeSubHeadingListStart}{\begin{itemize}}
\newcommand{\resumeSubHeadingListEnd}{\end{itemize}}
\newcommand{\resumeItemListStart}{\begin{itemize}}
\newcommand{\resumeItemListEnd}{\end{itemize}}
\newcommand{\projectSubheadingListStart}{\begin{itemize}}
\newcommand{\projectSubheadingListEnd}{\end{itemize}}

% -------------------- 9. SPACING HELPERS -----------------------------
\newcommand{\educationMulticolVSpaceBefore}{\vspace{-2pt}}
\newcommand{\educationMulticolVSpaceAfter}{\vspace{-11pt}}
\newcommand{\headerBottomVSpace}{\vspace{2.5pt}}

% ============================================================
%  END ENGINE
% ------------------------------------------------------------
%  COMMAND REFERENCE -- the full editable API for this template
%
%  \resumeSubheading{Org}{Location}{Role}{Dates}
%      Two-line header entry. Org/Location on row 1, Role/Dates on row 2.
%      Used by: Training & Development, Model Releases.
%
%  \projectSubheading{Title}{Subtitle}{Date}
%      One-line header entry. Title + italic subtitle on the left, date right.
%      Used by: Projects & Research, Industry Initiatives.
%
%  \cvCertification{Title}{Issuer}{Date}
%      Two-line header entry for certifications/courses/evaluations.
%      Used by: Safety, Evaluation & Compliance.
%
%  \resumeItem{text}                 -- a single bullet
%  \resumeSubItem{Bold label}{text}  -- a bullet with a bold leading label
%
%  \resumeSubHeadingListStart / \resumeSubHeadingListEnd
%      Wraps a block of \resumeSubheading entries.
%  \projectSubheadingListStart / \projectSubheadingListEnd
%      Wraps a block of \projectSubheading or \cvCertification entries.
%  \resumeItemListStart / \resumeItemListEnd
%      Wraps the bullets belonging to one entry.
%  \resumeSubBulletListStart / \resumeSubBulletListEnd
%      A dash-style nested bullet list, one level deeper than normal.
%  compactitemize (environment)
%      Bullet list used for Technical Capabilities / Languages rows.
%
%  Bullet length -- font (10pt Charter) and column width (0.97\textwidth)
%  are the drafting reference; no fixed character count (CV - Rules.md
%  Rule 1).
% ============================================================

\begin{document}

% ============================================================
%  CONTENT ZONE -- everything below is safe to edit/regenerate.
%  This is the ONE source-of-truth file: every section, every
%  entry, every bullet shape the engine supports. Compiled as-is,
%  it IS the full master CV, and the only file needed to derive a
%  1-page version -- no separate file required. Page count changes
%  whenever entries change; verify by compiling.
%
%  Every entry below carries a "1-PAGE TIER" tag (CORE/TRIM/CUT) --
%  see `CV - Rules.md` Rule 33 for what they mean and how to derive
%  a 1-page version from this file.
%
%  CONTENT DISCLAIMER: the "Claude" entries below exist to demonstrate
%  command usage and the 1-PAGE TIER system only. Wording is
%  illustrative, not a rule-compliance reference -- see
%  `CV - Rules.md` for the actual CV rules. Do not copy phrasing
%  patterns from this file into a real CV.
% ============================================================

% ---------------- HEADER ----------------
\begin{documentHeader}
\textbf{{\Huge Claude}}\\
\vspace{-8pt}
\Large Frontier AI Model Family \textemdash{} Anthropic, PBC

\vspace{-4pt}
\normalsize
\kern 1.0pt
\vspace{2pt}

\mbox{{\color{black}\footnotesize\faMapMarker*}\hspace*{0.13cm}San Francisco, California, USA}%
\kern 0.25 cm%
\ContactSeparator%
\kern 0.25 cm%
\mbox{\textcolor{black}{\footnotesize\faEnvelope[regular]} \href{mailto:claude@anthropic.com}{claude@anthropic.com}}%
\kern 0.25 cm%
\ContactSeparator%
\kern 0.25 cm%
\mbox{\textcolor{black}{\footnotesize\faPhone*} \href{tel:+14155550188}{+1 (415) 555-0188}}\\
\mbox{\textcolor{black}{\footnotesize\faLinkedinIn} \href{https://www.linkedin.com/company/anthropicresearch/}{anthropicresearch}}%
\kern 0.25 cm%
\ContactSeparator%
\kern 0.25 cm%
\mbox{\textcolor{black}{\footnotesize\faGithub} \href{https://github.com/anthropics}{anthropics}}
\headerBottomVSpace
\end{documentHeader}
% NOTE: email/phone above are illustrative placeholders (Claude has no
% personal inbox or phone line) -- everything else is a real, live link.
% Header is always kept -- there is no 1-PAGE TIER tag for it.

% ---------------- TRAINING & DEVELOPMENT ----------------
\section{Training \& Development}
\resumeSubHeadingListStart
% [1-PAGE TIER: TRIM -- keep the heading + multicols stat pair; cut the
%  fine-tuning-cadence sub-entry and its dash sub-bullet below.]
% NOTE: full Constitutional AI facts (RLAIF, written constitution) live
% once, under Projects & Research below. This sub-entry demonstrates
% \resumeSubBulletListStart with a DIFFERENT, non-duplicated fact --
% don't restate the Projects & Research content here. See `CV - Rules.md`,
% "On duplication."
\resumeSubheading
  {Anthropic, PBC}{San Francisco, CA}
  {Pretraining, Fine-Tuning \& Constitutional AI}{Jan 2021 -- Present}
\resumeItemListStart
\educationMulticolVSpaceBefore
\begin{multicols}{2}
  \resumeItem{\textbf{Knowledge cutoff:} May 2025.}
  \resumeItem{\textbf{Context window:} up to \textbf{1M} tokens.}
\end{multicols}
\educationMulticolVSpaceAfter
\resumeItem{\textbf{Fine-Tuning Cadence} \textit{-- Post-Training} \hfill Ongoing}
\resumeSubBulletListStart
  \resumeItem{Iterates through \textbf{RLHF} and targeted fine-tuning passes between major model releases.}
\resumeSubBulletListEnd
\resumeItemListEnd
\resumeSubHeadingListEnd

% ---------------- MODEL RELEASES ----------------
\section{Model Releases}
\resumeSubHeadingListStart

% [1-PAGE TIER: TRIM -- keep bullets 1-3, cut bullet 4]
% Claude Opus 4.8 / Sonnet 5 / Haiku 4.5
\resumeSubheading
  {Anthropic, PBC \textmd{\textit{-- \href{https://www.anthropic.com/news}{Announcements}}}}{San Francisco, CA}
  {\textbf{Claude Opus 4.8 / Sonnet 5 / Haiku 4.5} -- Flagship Model Family}{May 2026 -- Present}
\resumeItemListStart
  \resumeItem{Serves as Anthropic's flagship lineup across \textbf{Opus}, \textbf{Sonnet} \& \textbf{Haiku}, powering Claude.ai.}
  \resumeItem{Extends \textbf{Claude 4}-era agentic workflows across \textbf{Claude Code}, \textbf{Claude in Chrome} \& \textbf{Cowork}.}
  \resumeItem{Operates under Anthropic's \textbf{Responsible Scaling Policy}, vetted by independent evaluators.}
  \resumeItem{Supports a \textbf{200K}-token context window, with extended windows in select deployments.}
\resumeItemListEnd

% [1-PAGE TIER: TRIM -- keep bullets 1-3, cut bullet 4]
% Claude 4 Generation
\resumeSubheading
  {Anthropic, PBC}{San Francisco, CA}
  {\textbf{Claude 4 Generation} -- Opus 4 \& Sonnet 4}{May 2025 -- Apr 2026}
\resumeItemListStart
  \resumeItem{Brought \textbf{Claude Code} to general availability across terminal, IDE \& web.}
  \resumeItem{Introduced extended, interleaved thinking with \textbf{parallel tool-calling} for long tasks.}
  \resumeItem{Established \textbf{computer use} as a standard agentic capability across the model family.}
  \resumeItem{Expanded enterprise adoption via \textbf{Claude for Enterprise} and deeper \textbf{MCP} integrations.}
\resumeItemListEnd

% [1-PAGE TIER: CORE -- all 3 bullets now fit; keep as-is]
% Claude 3.7 Sonnet
\resumeSubheading
  {Anthropic, PBC}{San Francisco, CA}
  {\textbf{Claude 3.7 Sonnet -- Hybrid Reasoning Model} -- \href{https://www.anthropic.com/news}{Announcement}}{Feb 2025 -- May 2025}
\resumeItemListStart
  \resumeItem{Introduced \textbf{hybrid reasoning}: a standard mode plus an extended thinking mode.}
  \resumeItem{Launched alongside \textbf{Claude Code}, Anthropic's first agentic command-line coding tool.}
  \resumeItem{Improved agentic coding, becoming a preferred model for \textbf{SWE-bench}-style tasks.}
\resumeItemListEnd

% [1-PAGE TIER: CUT -- omit this entire entry for a 1-page version]
% Claude 3.5 Sonnet (New) -- Computer Use
\resumeSubheading
  {Anthropic, PBC \textmd{\textit{-- \href{https://www.anthropic.com/news/3-5-models-and-computer-use}{Announcement}}}}{San Francisco, CA}
  {\textbf{Claude 3.5 Sonnet (New)} -- Computer Use Pioneer}{Oct 2024 -- Feb 2025}
\resumeItemListStart
  \resumeItem{Shipped \textbf{computer use} in beta -- the first model able to view a screen and act on it.}
  \resumeItem{Upgraded coding performance over the original \textbf{3.5 Sonnet}, same \textbf{200K} context.}
  \resumeItem{Opened computer use to developers via the \textbf{API}, enabling desktop-control agents.}
\resumeItemListEnd

% [1-PAGE TIER: CUT]
% Claude 3.5 Sonnet
\resumeSubheading
  {Anthropic, PBC}{San Francisco, CA}
  {\textbf{Claude 3.5 Sonnet}}{Jun 2024 -- Oct 2024}
\resumeItemListStart
  \resumeItem{Outperformed the larger \textbf{Claude 3 Opus} on benchmarks at roughly \textbf{80\%} lower cost.}
  \resumeItem{Introduced \textbf{Artifacts}, a side-by-side workspace for code, documents \& diagrams.}
  \resumeItem{Became Anthropic's fastest mid-tier model at the time, balancing intelligence with low latency.}
\resumeItemListEnd

% [1-PAGE TIER: CUT]
% Claude 3 Family
\resumeSubheading
  {Anthropic, PBC}{San Francisco, CA}
  {\textbf{Claude 3 Family} -- Opus, Sonnet \& Haiku}{Mar 2024 -- Jun 2024}
\resumeItemListStart
  \resumeItem{Launched as a \textbf{3}-tier family (\textbf{Opus}, \textbf{Sonnet}, \textbf{Haiku}).}
  \resumeItem{Added native \textbf{vision} support, reading charts \& diagrams directly from images.}
  \resumeItem{Extended context window to \textbf{200K} tokens across the family at general availability.}
\resumeItemListEnd

% [1-PAGE TIER: CUT]
% Claude 1 / Claude Instant
\resumeSubheading
  {Anthropic, PBC}{San Francisco, CA}
  {\textbf{Claude} \& \textbf{Claude Instant}}{Mar 2023 -- Mar 2024}
\resumeItemListStart
  \resumeItem{Shipped as Anthropic's first publicly available models via the Claude API.}
  \resumeItem{Established \textbf{Constitutional AI} as the core alignment method from day one.}
  \resumeItem{Reached a \textbf{100K}-token context window with \textbf{Claude 2.1}, then among the largest available.}
\resumeItemListEnd

\resumeSubHeadingListEnd

% ---------------- PROJECTS & RESEARCH ----------------
% Ordered reverse-chronologically by END date (Present/Ongoing = most
% recent), matching Rule 10 in `CV - Rules.md` -- not by start date.
% Where two entries are both open-ended, the one with an actual dated
% start (e.g. "Feb 2025 -- Present") is placed above a bare "Ongoing"
% with no start date given, as the more specifically recent of the two.
\section{Projects \& Research}
\projectSubheadingListStart

% [1-PAGE TIER: CUT]
% Claude Code
\projectSubheading
  {Claude Code}
  {-- \href{https://github.com/anthropics/claude-code}{GitHub} -- \href{https://docs.claude.com}{Docs}}
  {Feb 2025 -- Present}
\resumeItemListStart
  \resumeItem{Brought agentic coding to the terminal, IDE \& web, reading \& editing full repositories.}
  \resumeItem{Reached general availability alongside \textbf{Claude 4}, then added a web interface.}
\resumeItemListEnd

% [1-PAGE TIER: CUT]
% Anthropic Cookbook
\projectSubheading
  {Anthropic Cookbook}
  {-- \href{https://github.com/anthropics/anthropic-cookbook}{GitHub}}
  {Ongoing}
\resumeItemListStart
  \resumeItem{Maintains open-source, runnable examples for retrieval, tool use \& multimodal prompting.}
  \resumeItem{Used by developers as the reference starting point for building on the \textbf{Claude API}.}
\resumeItemListEnd

% [1-PAGE TIER: TRIM -- keep bullets 1 \& 3, cut bullet 2]
% Model Context Protocol
\projectSubheading
  {Model Context Protocol}
  {-- \href{https://www.anthropic.com/news/model-context-protocol}{Announcement} -- \href{https://github.com/modelcontextprotocol}{GitHub}}
  {Nov 2024}
\resumeItemListStart
  \resumeItem{Open-sourced an \textbf{open standard} connecting models to tools, files \& live data.}
  \resumeItem{Shipped SDKs in \textbf{Python} \& \textbf{TypeScript} plus reference servers.}
  \resumeItem{Adopted industry-wide as a common protocol for agentic tool use beyond Anthropic's own models.}
\resumeItemListEnd

% [1-PAGE TIER: CUT]
% Scaling Monosemanticity
\projectSubheading
  {Scaling Monosemanticity}
  {-- \href{https://transformer-circuits.pub/2024/scaling-monosemanticity}{Paper}}
  {May 2024}
\resumeItemListStart
  \resumeItem{Extended sparse-autoencoder methods to \textbf{Claude 3 Sonnet}, a production-scale model.}
  \resumeItem{Surfaced millions of interpretable features, including deception \& bias concepts.}
\resumeItemListEnd

% [1-PAGE TIER: TRIM -- keep bullets 1 \& 3, cut bullet 2]
% Constitutional AI paper
\projectSubheading
  {Constitutional AI}
  {-- \href{https://arxiv.org/abs/2212.08073}{Paper}}
  {Dec 2022}
\resumeItemListStart
  \resumeItem{Proposed training a model to critique \& revise outputs against a written constitution.}
  \resumeItem{Reduced reliance on human labeling for harmlessness, using \textbf{AI feedback} (RLAIF).}
  \resumeItem{Became the foundation alignment method for every subsequent Claude release.}
\resumeItemListEnd

% NOTE: third-party pre-deployment evaluation (METR, Apollo Research,
% RSP/ASL thresholds) lives once, under Safety, Evaluation & Compliance
% below -- not repeated here as a separate Projects & Research entry.

\projectSubheadingListEnd

% ---------------- TECHNICAL CAPABILITIES ----------------
% [1-PAGE TIER: TRIM -- keep all 5 rows, but tighten wording if tight on
%  space, e.g. drop "code execution", "/JavaScript", "RSP compliance".]
\section{Technical Capabilities}
\begin{compactitemize}
  \resumeSubItem{\textbf{Reasoning \& Tool Use:}}{Extended thinking, tool-calling, computer use, code execution.}
  \resumeSubItem{\textbf{Software Engineering:}}{Python, TypeScript, Bash, Git, Claude Code, Claude Agent SDK.}
  \resumeSubItem{\textbf{Multimodal Understanding:}}{Vision (images, screenshots, PDFs), document \& diagram parsing.}
  \resumeSubItem{\textbf{Integrations \& Protocols:}}{Model Context Protocol (MCP), function calling, RAG, REST/SDKs.}
  \resumeSubItem{\textbf{Safety \& Alignment Techniques:}}{Constitutional AI, RLHF, red-teaming, interpretability.}
\end{compactitemize}

% ---------------- SAFETY, EVALUATION & COMPLIANCE ----------------
\section{Safety, Evaluation \& Compliance}
\projectSubheadingListStart

% [1-PAGE TIER: TRIM -- keep bullet 1 only, cut bullet 2]
% Third-Party Pre-Deployment Evaluation
\cvCertification
  {Third-Party Pre-Deployment Evaluation}
  {METR, Apollo Research \& UK / US AI Safety Institutes}{Ongoing}
\resumeItemListStart
  \resumeItem{Undergoes independent testing for dangerous-capability \& autonomy risks pre-release.}
  \resumeItem{Findings feed directly into the published \textbf{model/system card} for each release.}
\resumeItemListEnd

% [1-PAGE TIER: CUT]
% Bug Bounty & Responsible Disclosure
\cvCertification
  {Bug Bounty \& Responsible Disclosure}
  {Anthropic, PBC -- \href{https://www.anthropic.com}{HackerOne Program}}{Ongoing}
\resumeItemListStart
  \resumeItem{Runs a public vulnerability disclosure program covering the Claude API \& consumer products.}
  \resumeItem{Rewards researchers for responsibly reporting jailbreaks, safety bypasses \& security issues.}
\resumeItemListEnd

\projectSubheadingListEnd

% ---------------- INDUSTRY INITIATIVES ----------------
% [1-PAGE TIER: CUT -- omit this entire section for a 1-page version]
\section{Industry Initiatives}
\projectSubheadingListStart
\projectSubheading
  {Frontier Model Forum -- Founding Member}
  {}
  {Jul 2023 -- Present}
\resumeItemListStart
  \resumeItem{Co-founded with \textbf{OpenAI}, \textbf{Google DeepMind} \& \textbf{Microsoft} on safety research.}
  \resumeItem{Shares best practices on evaluation, red-teaming \& responsible scaling across member labs.}
  \resumeItem{Coordinates with policymakers \& researchers on safety standards for frontier AI systems.}
\resumeItemListEnd
\projectSubheadingListEnd

% ---------------- LANGUAGES ----------------
% [1-PAGE TIER: CORE -- always keep, unchanged]
\section{Languages}
\vspace{-12pt}
\begin{compactitemize}
\begin{multicols}{2}
  \resumeSubItem{\textbf{English:}}{Native fluency.}
  \resumeSubItem{\textbf{Spanish:}}{Fluent.}
  \resumeSubItem{\textbf{Mandarin:}}{Fluent.}
  \resumeSubItem{\textbf{Japanese:}}{Fluent.}
\end{multicols}
\end{compactitemize}

\end{document}
